\section{Źródła i podstawa opracowania}
\small
Źródła otwarte sprawdzono 19.10.2026. Opisy katalogowe norm potwierdzają ich temat i oznaczenia; nie zastępują dostępu do pełnej treści. Projekt nie deklaruje sprawdzenia wszystkich postanowień norm ani uzyskania wymaganych uzgodnień. Obliczenia, geometria, scenariusze i wymagania liczbowe są własnymi założeniami dydaktycznymi.

\begin{enumerate}[label={[\arabic*]},itemsep=8pt,leftmargin=8mm]
\item Rządowe Centrum Bezpieczeństwa. \textit{Infrastruktura krytyczna}. Definicja i rozróżnienie między obiektem strategicznym a formalną IK.\\\url{https://www.gov.pl/web/rcb/infrastruktura-krytyczna}
\item MSWiA. \textit{Ochrona przeciwpożarowa budynków, innych obiektów budowlanych i terenów}. Tekst jednolity: Dz.U. 2023 poz. 822, wraz z późniejszymi zmianami wskazanymi w ELI. Przed realizacją sprawdzić aktualny stan prawny.\\\url{https://eli.gov.pl/eli/DU/2023/822/ogl}
\item ONVIF. \textit{Profile T}. Funkcje transmisji wideo i interoperacyjność; zgodność konkretnego modelu wymaga sprawdzenia.\\\url{https://www.onvif.org/profiles/profile-t/}
\item Security Industry Association. \textit{Open Supervised Device Protocol (OSDP)}. Komunikacja czytnik--kontroler i Secure Channel.\\\href{https://www.securityindustry.org/industry-standards/open-supervised-device-protocol/}{SIA -- specyfikacja i opis OSDP}
\item Europejska Rada Ochrony Danych. \textit{Guidelines 3/2019 on processing of personal data through video devices}, wersja 2.1. Cel, retencja, przejrzystość, środki ochrony i ocena skutków.\\\href{https://www.edpb.europa.eu/sites/default/files/files/file1/edpb_guidelines_201903_video_devices_en_0.pdf}{Dokument EROD -- wytyczne dotyczące monitoringu wizyjnego (PDF)}
\item Polski Komitet Normalizacyjny. \textit{Komunikat on-line}. PN-EN IEC 62676-4:2026-01 wskazana jako zastępująca wydanie PN-EN 62676-4:2015-06; zakres: wytyczne stosowania CCTV.\\\url{https://www.pkn.pl/komunikat-online}
\item Polski Komitet Normalizacyjny. \textit{PN-EN 50131-1:2009 -- Wymagania systemowe SSWiN}, wraz z właściwymi zmianami. Opis katalogowy ogranicza zakres do systemów w budynkach.\\\url{https://sklep.pkn.pl/pn-en-50131-1-2009p.html}
\item Polski Komitet Normalizacyjny. \textit{PN-EN 54-4:2001 -- Zasilacze}, wraz z A1 i A2. Punkt odniesienia do doboru zasilaczy SSP; inne części serii EN 54 dobiera się do komponentów.\\\url{https://sklep.pkn.pl/pn-en-54-4-2001p.html}
\item CISA / NSA. \textit{Top Ten Cybersecurity Misconfigurations}, AA23-278A. Segmentacja, kontrola dostępu i ograniczanie połączeń.\\\href{https://www.cisa.gov/news-events/cybersecurity-advisories/aa23-278a}{CISA -- zalecenia dotyczące konfiguracji i segmentacji}
\item NIST. \textit{AI Risk Management Framework}. Odniesienie do oceny ryzyka i walidacji analityki.\\\url{https://www.nist.gov/itl/ai-risk-management-framework}
\end{enumerate}

\vfill
\begin{note}{Pochodzenie ilustracji}
Okładka: obraz wygenerowany wbudowanym narzędziem ImageGen; pełny prompt w repozytorium. Rysunki techniczne: własne grafiki wektorowe generowane ze źródeł Python/ReportLab. Nie są planami żadnego rzeczywistego lotniska. AI wspomogła przygotowanie tekstu i kodu; zespół powinien zweryfikować i przyjąć treść przed oddaniem.
\end{note}
